\documentclass[11pt]{article}

\usepackage[a4paper,margin=1in]{geometry}
\usepackage{amsmath,amssymb,mathtools}
\usepackage[hidelinks]{hyperref}

\title{Implementation Summary}
\author{}
\date{}

\begin{document}
\maketitle

This file records what was implemented from the two-agent plan and what remains
as the next mathematical work.

\section{New Artifacts}

\begin{enumerate}
\item \texttt{level-set-deficit-identity.md}

This implements the Nicola--Tilli route up to an exact theorem:
\[
\frac{2}{|\Omega|}\big(E(\Omega)-E(B)\big)
=
\frac1{|\Omega|}
\int_0^{\|u\|_\infty}
\frac{D_H(t)+D_I(t)}
{n^2\omega_n^{2/n}m(t)^{1-2/n}}\,dt.
\]
It also proves the weighted variance identity
\[
\int_{\Sigma_t}
\frac{(|\nabla u|-\bar f)^2}{|\nabla u|}
=
\frac{m(t)}{P(t)^2}D_H(t).
\]
This gives a precise bridge from the convexity gap in
\texttt{faber-krahn-1.pdf} to approximate Serrin data on good superlevel sets.

\item \texttt{quantitative-selection-principle.md}

This implements the BDV selection route as a single-set lemma. Given
\[
\varepsilon=\alpha(\Omega_0),\qquad
\delta=E(\Omega_0)-E(B_1),\qquad
q=\delta/\varepsilon\ll1,
\]
it selects a minimizer of
\[
E(U)+f_{\hat\eta}(|U|)
+\sqrt{\varepsilon^2+\tau^2(\alpha(U)-\varepsilon)^2},
\qquad
\tau=q^{1/4},
\]
and proves the quantitative preservation estimates
\[
\alpha(\widetilde\Omega)
\ge
(1-Cq^{1/4})\alpha(\Omega_0),
\]
\[
E(\widetilde\Omega)-E(B_1)
\le
C(E(\Omega_0)-E(B_1)).
\]
Hence
\[
Q_\alpha(\widetilde\Omega)
\le
C Q_\alpha(\Omega_0).
\]
\end{enumerate}

\section{What Is Now Proved or Reduced}

The Talenti/Nicola--Tilli route now has a clean first theorem: the global
Saint-Venant deficit is exactly the integrated level-set loss. This is a
standalone quantitative refinement of the proof in \texttt{faber-krahn-1.pdf}.

The selection-principle route now has a concrete replacement for the sequence
argument in BDV Proposition 4.4: the selected set preserves the relevant
deficit-to-\(\alpha\) ratio up to constants.

The compactness bottleneck has been localized. The remaining non-explicit part
is not the penalized minimization; it is the quantitative entry into the global
nearly spherical graph regime.

\section{Next Targets}

\begin{enumerate}
\item Write the finite-perimeter version of the level-set identity.

Replace smooth regular-level arguments by a.e. coarea statements for the
torsion function. This should be mostly technical and self-contained.

\item Prove a boundary-layer propagation lemma.

The target form is
\[
\exists t\downarrow0
\quad\hbox{with}\quad
|\Omega\setminus\{u>t\}|\hbox{ small}
\quad\hbox{and}\quad
D_H(t)+D_I(t)\hbox{ controlled}.
\]
This is the main missing step for making the level-set route prove sharp
Fraenkel stability of the original domain.

\item Track the Alt-Caffarelli flatness constants for selected minimizers.

The needed statement is
\[
\alpha(U_*)\le\varepsilon_0(N,R,\tau)
\Rightarrow
\partial U_*
\hbox{ is a global }C^{1,\gamma}\hbox{ graph over }\partial B_1.
\]
The new note already proves the preceding Hausdorff-closeness estimate from
density and \(\alpha\).

\item Collect stable Serrin inputs compatible with the level-set variance.

Existing stable Serrin theorems often use oscillation or Lipschitz norms of the
normal derivative. The level-set route naturally gives a weighted \(L^2\)
variance of \(|\nabla u|\), so the next search should focus on Serrin stability
results accepting integral boundary data.
\end{enumerate}

\end{document}
